%------------------------------------------------------------------------------%
\documentclass[a4paper,12pt,fleqn]{article}

\usepackage{calculo_varias_variaveis-1}

\ShowAnswers

%------------------------------------------------------------------------------%
\begin{document}

\makeHeader{CFVVI}{Prova 4 -- Turma 2}{24/11/2025}

% 01
%------------------------------------------------------------------------------%
\exercise{50\%}
Determine e classifique os pontos críticos da função
\(
  f(x,y) = \ln(1+x^2+y^2) - xy
\)

\clearpagequestiononly

\begin{answer}
  Calculando o gradiente
  \[
    f_x(x, y)
    = \D{}{x}\left(\ln(1+x^2+y^2) - xy\right)
    = \frac{2x}{1+x^2+y^2} - y
  \]
  \[
    f_y(x, y)
    = \D{}{y}\left(\ln(1+x^2+y^2) - xy\right)
    = \frac{2y}{1+x^2+y^2} - x
  \]
  Igualando o gradiente ao vetor nulo
  \[
    \begin{cases}
      \dfrac{2x}{1+x^2+y^2} - y = 0 \\[4mm]
      \dfrac{2y}{1+x^2+y^2} - x = 0
    \end{cases}
  \]
  rearranjando
  \[
    \begin{cases}
      2x = y\left(1+x^2+y^2\right) \\
      2y = x\left(1+x^2+y^2\right)
    \end{cases}
  \]
  Subtraindo as equações
  \[
    2(x-y) = -(x-y)\left(1+x^2+y^2\right)
  \]

  Caso 1: $x\neq y$, dividimos os dois lados por $x-y$, obtendo
  \begin{align*}
    2(x-y)    & = -(x-y)\left(1+x^2+y^2\right) \\
    2         & = -\left(1 + x^2 + y^2\right)  \\
    2         & = -1 - x^2 - y^2               \\
    x^2 + y^2 & = -3
  \end{align*}
  Não existe solução

  \bigskip
  Caso 2: $x=y$, a equação \(2(x-y) = -(x-y)\left(1+x^2+y^2\right)\) é satisfeita

  Substituindo na primeira equação e obtemos
  \begin{align*}
    2x = y\left(1+x^2+y^2\right) \\
    2x = x\left(1+x^2+x^2\right) \\
    2x = x\left(1+2x^2\right)
  \end{align*}

  Se $x=0$, essa equação é satisfeita e temos o primeiro ponto crítico \(P_1=(0, 0)\)

  Se $x\neq y$, dividimos os dois termos por $x$
  \begin{align*}
    2x      & = x\left(1+2x^2\right) \\
    2       & = 1 + 2x^2             \\
    2x^2    & = 1                    \\
    x^2     & = \frac{1}{2}          \\
    \abs{x} & = \frac{1}{\sqrt{2}}   \\
    x       & = \frac{\pm 1}{\sqrt{2}}
  \end{align*}
  Encontramos os pontos críticos
  \(P_2 = \left( \frac{1}{\sqrt{2}},  \frac{1}{\sqrt{2}}\right)\) e
  \(P_3 = \left(-\frac{1}{\sqrt{2}}, -\frac{1}{\sqrt{2}}\right)\)

  \bigskip
  Calculando a Hessiana de $f$
  \[
    f_{xx}(x, y)
    = \D{}{x}\left(\frac{2x}{1+x^2+y^2} - y\right)
    = \frac{2(1+x^2+y^2)-4x^2}{\left(1+x^2+y^2\right)^2}
    = \frac{2-2x^2+2y^2}{\left(1+x^2+y^2\right)^2}
  \]
  \[
    f_{yy}(x, y)
    = \D{}{y}\left(\frac{2y}{1+x^2+y^2} - x\right)
    = \frac{2(1+x^2+y^2)-4y^2}{\left(1+x^2+y^2\right)^2}
    = \frac{2+2x^2-2y^2}{\left(1+x^2+y^2\right)^2}
  \]
  \[
    f_{xy}(x, y)
    = \D{}{x}\left(\frac{2y}{1+x^2+y^2} - x\right)
    = \frac{-4xy}{\left(1+x^2+y^2\right)^2} -1
  \]
  Discriminante (Determinante da Hessiana)
  \[
    D(x, y) = f_{xx}(x, y)f_{yy}(x, y) - f^2_{xy}(x, y)
  \]

  \bigskip
  Classificando o ponto \((0, 0)\)
  \[
    f_{xx}(0, 0)
    = \left(\frac{2-2x^2+2y^2}{\left(1+x^2+y^2\right)^2}\right)\evalat{(0, 0)}{}
    = 2
  \]
  \[
    f_{yy}(0, 0)
    = \left(\frac{2+2x^2-2y^2}{\left(1+x^2+y^2\right)^2}\right)\evalat{(0, 0)}{}
    = 2
  \]
  \[
    f_{xy}(0, 0)
    = \left(\frac{-4xy}{\left(1+x^2+y^2\right)^2} -1\right)\evalat{(0, 0)}{}
    = -1
  \]
  \[
    D(0, 0)
    = 2 \times 2 - (-1)^2
    = 3 > 0
  \]
  Como \(f_{xx}(0, 0) = 2 > 0\), é um \textbf{mínimo local}

  \bigskip
  No ponto \(\left( \frac{1}{\sqrt{2}},  \frac{1}{\sqrt{2}}\right)\)
  temos
  \[
    \left(\left(1+x^2+y^2\right)^2\right)\evalat{\left( \frac{1}{\sqrt{2}},  \frac{1}{\sqrt{2}}\right)}{}
    = \left(1+\frac{1}{2}+\frac{1}{2}\right)^2
    = 2^2
    = 4
  \]
  \[
    f_{xx}\left( \frac{1}{\sqrt{2}},  \frac{1}{\sqrt{2}}\right)
    = \left(\frac{2-2x^2+2y^2}{\left(1+x^2+y^2\right)^2}\right)\evalat{\left( \frac{1}{\sqrt{2}},  \frac{1}{\sqrt{2}}\right)}{}
    = \frac{2-1+1}{4}
    = \frac{1}{2}
  \]
  \[
    f_{yy}\left( \frac{1}{\sqrt{2}},  \frac{1}{\sqrt{2}}\right)
    = \left(\frac{2+2x^2-2y^2}{\left(1+x^2+y^2\right)^2}\right)\evalat{\left( \frac{1}{\sqrt{2}},  \frac{1}{\sqrt{2}}\right)}{}
    = \frac{2+1-1}{4}
    = \frac{1}{2}
  \]
  \[
    f_{xy}\left( \frac{1}{\sqrt{2}},  \frac{1}{\sqrt{2}}\right)
    = \left(\frac{-4xy}{\left(1+x^2+y^2\right)^2} -1\right)\evalat{\left( \frac{1}{\sqrt{2}},  \frac{1}{\sqrt{2}}\right)}{}
    = \frac{-2}{4} -1
    = \frac{-3}{2}
  \]
  \[
    D\left( \frac{1}{\sqrt{2}},  \frac{1}{\sqrt{2}}\right)
    = \frac{1}{2}\times\frac{1}{2} - \left(\frac{-3}{2}\right)^2
    = \frac{1}{4} - \frac{9}{4}
    = \frac{-8}{4}
    = -2
    < 0
  \]
  Portanto, é um \textbf{ponto de sela}

  \bigskip
  No ponto \(\left( -\frac{1}{\sqrt{2}}, -\frac{1}{\sqrt{2}}\right)\)
  temos
  \[
    \left(\left(1+x^2+y^2\right)^2\right)\evalat{\left(-\frac{1}{\sqrt{2}}, -\frac{1}{\sqrt{2}}\right)}{}
    = \left(1+\frac{1}{2}+\frac{1}{2}\right)^2
    = 2^2
    = 4
  \]
  \[
  f_{xx}\left( \frac{1}{\sqrt{2}},  \frac{1}{\sqrt{2}}\right)
  = \left(\frac{2-2x^2+2y^2}{\left(1+x^2+y^2\right)^2}\right)\evalat{\left(-\frac{1}{\sqrt{2}}, -\frac{1}{\sqrt{2}}\right)}{}
  = \frac{2-1+1}{4}
  = \frac{1}{2}
  \]
  \[
  f_{yy}\left( \frac{1}{\sqrt{2}},  \frac{1}{\sqrt{2}}\right)
  = \left(\frac{2+2x^2-2y^2}{\left(1+x^2+y^2\right)^2}\right)\evalat{\left(-\frac{1}{\sqrt{2}}, -\frac{1}{\sqrt{2}}\right)}{}
  = \frac{2+1-1}{4}
  = \frac{1}{2}
  \]
  \[
  f_{xy}\left( \frac{1}{\sqrt{2}},  \frac{1}{\sqrt{2}}\right)
  = \left(\frac{-4xy}{\left(1+x^2+y^2\right)^2} -1\right)\evalat{\left(-\frac{1}{\sqrt{2}}, -\frac{1}{\sqrt{2}}\right)}{}
  = \frac{-2}{4} -1
  = \frac{-3}{2}
  \]
  \[
  D\left(-\frac{1}{\sqrt{2}}, -\frac{1}{\sqrt{2}}\right)
  = \frac{1}{2}\times\frac{1}{2} - \left(\frac{-3}{2}\right)^2
  = \frac{1}{4} - \frac{9}{4}
  = \frac{-8}{4}
  = -2
  < 0
  \]
  Portanto, é um \textbf{ponto de sela}
\end{answer}

% 02
%------------------------------------------------------------------------------%
\exercise{50\%}
Determine os valores máximo e mínimo absolutos da função
\(
  f(x, y) = x^2 + y^2 - 2x - 2y
\)
na região fechada imitada
\(
  R = \{(x, y)\in\R^2\colon \abs{x}\le 2 ,\; \abs{y}\le 4 \}
\)

\clearpagequestiononly

\begin{answer}
\textbf{Interior}

Derivadas parciais
\[
  f_x(x, y) = 2x - 2 \qquad
  f_y(x, y) = 2y - 2
\]
Igualando a zero
\[
  2x-2=0\quad\Rightarrow\quad x=1
\]
\[
  2y-2=0\quad\Rightarrow\quad y=1
\]
Ponto crítico é $(1, 1)$, como \(\abs{1}\le 2\) e \(\abs{1}\le 4\),
o ponto crítico está no interior de $R$

Valor da função no ponto crítico
\[
  f(1, 1)
  = \left(x^2 + y^2 - 2x - 2y\right) \evalat{(1, 1)}{}
  = 1 + 1 - 2 - 2
  = -2
\]

\bigskip
A fronteira de \(R\) é composta por quatro arestas
\begin{align*}
  x = -2     & \quad y\in[-4, 4] \\
  x = \hpm 2 & \quad y\in[-4, 4] \\
  y = - 4    & \quad x\in[-2, 2] \\
  y = \hpm 4 & \quad x\in[-2, 2]
\end{align*}


\bigskip
\textbf{Aresta \(x=2\)}
\[
  g(y)
  = f(2, y)
  = 4 + y^2 - 4 - 2y
  = y^2 - 2y
\]
\[
  g'(y) = 2y - 2
\]
cuja raiz é \(y=1\)

Então os candidatos a extremos no problema 1D são $y=-4$, $y=1$ e $y=4$

Os candidatos no problema 2D são $(2, -4)$, $(2, 1)$ e $(2, 4)$

Avaliando $f$
\[
  f(2, -4)
  = g(-4)
  = 16+8
  = 24
\]
\[
  f(2, 1)
  = g(1)
  = -1
\]
\[
  f(2, 4)
  = g(4)
  = 16-8
  = 8
\]


\bigskip
\textbf{Aresta \(x=-2\)}
\[
  h(y)
  = f(-2,y)
  = 4+y^{2}+4-2y
  = y^{2}-2y+8
\]
\[
  h'(y)=2y-2
\]
com raiz \(y=1\)

Então os candidatos a extremos no problema 1D são $y=-4$, $y=1$ e $y=4$

Os candidatos no problema 2D são $(-2, -4)$, $(-2, 1)$ e $(-2, 4)$

Avaliando $f$
\[
  f(-2, -4)
  = h(-4)
  = 16 + 8 + 8
  = 32
\]
\[
  f(-2, 1)
  = h(1)
  = 1 - 2 + 8
  = 7
\]
\[
  f(-2, 4)
  = h(4)
  = 16 - 8 + 8
  = 16
\]


\bigskip
\textbf{Aresta \(y=4\)}
\[
  p(x)
  = f(x, 4)
  = x^{2}+16-2x-8
  = x^{2}-2x+8
\]
\[
  p'(x)=2x-2
\]
com raiz \(x=1\)

Então os candidatos a extremos no problema 1D são $x=-2$, $x=1$ e $x=2$

Os candidatos no problema 2D são $(-2, 4)$, $(1, 4)$ e $(2, 4)$

Avaliando $f$
\[
  f(-2, 4)
  = p(-2)
  = 4+4+8
  = 16
\]
\[
  f(1, 4)
  = p(1)
  = 1-2+8
  = 7
\]
\[
  f(2, 4)
  = p(2)
  = 4-4+8
  = 8
\]

\bigskip
\textbf{Aresta \(y=-4\)}
\[
  q(x)
  = f(x, -4)
  = x^{2}+16-2x+8
  = x^{2}-2x+24
\]
\[
  q'(x) = 2x-2
\]
com raiz \(x=1\)

Então os candidatos a extremos no problema 1D são $x=-2$, $x=1$ e $x=2$

Os candidatos no problema 2D são $(-2, -4)$, $(1, -4)$ e $(2, -4)$

Avaliando $f$
\[
  f(-2, -4)
  = q(-2)
  = 4+4+24
  = 32
\]
\[
  f(1, -4)
  = q(1)
  = 1-2+24
  = 23
\]
\[
  f(2, -4)
  = q(2)
  = 4-4+24
  = 24
\]


\bigskip
\textbf{Comparando}
\begin{align*}
  f(1, 1)   & = -2 \\
  f(2, 4)   & = 8  \\
  f(2, -4)  & = 24 \\
  f(-2, 4)  & = 16 \\
  f(-2, -4) & = 32
\end{align*}

O mínimo absoluto de $f$ é $-2$ em $(1, 1)$

O máximo absoluto de $f$ é 32 em $(-2, -4)$
\end{answer}



%------------------------------------------------------------------------------%
\end{document}
%------------------------------------------------------------------------------%
